\section{Autonomous control}
\label{sec:control}

\subsection{The run: climb, decelerate, hold}
\label{sec:run}

An autonomous run has three stages (Fig.~\ref{fig:states}). In Climb, both motors run at full power. The robot has no way to measure distance, so Climb ends as soon as either of two conditions is met: \qty{3.15}{s} have passed since the start, or the pitch has stayed below \ang{18} for \qty{40}{ms}. The robot ignores the pitch condition for the first \qty{0.3}{s} of the run. In Decelerate, the drive power ramps linearly from full to zero over a set time, and momentum carries the robot from the ridge onto the mesa. It then enters Hold.

In Hold, the robot sets its motor power to zero while it is level, and the driver code applies zero power as short brake. Short brake connects each motor's two terminals together. A motor turned by an outside force then acts as a generator driving current through its own windings, and that current makes a torque opposing the motion, $\tau_\mathrm{brake} = k_\mathrm{t}^2 \omega / R_\mathrm{w}$ for a motor turned at speed $\omega$.

Hold also corrects the robot's position from its pitch. If the nose rises more than \ang{6}, because the robot stopped short of the top or is being shoved back down our ramp, it drives forward at full power, 255 of 255. If the nose drops more than \ang{6}, because the robot is about to go over the far side, it backs up at 110 of 255. It keeps correcting until the pitch is back within \ang{3}, so that it does not switch rapidly between correcting and braking at the edge of the \ang{6} band.

\begin{figure}[htbp]
  \centering
  \begin{tikzpicture}[
      font=\small,
      >={Stealth[length=2.4mm]},
      every edge/.style={draw,->,thick},
      state/.style={draw,thick,rounded corners=3pt,align=center,
                    text width=3.3cm,minimum height=1.35cm,inner sep=3pt},
      wait/.style={state,fill=black!6},
      run/.style={state,fill=blue!12},
      lbl/.style={align=center,font=\footnotesize,inner sep=2pt},
    ]
    \node[wait] (armed) at (0,0)       {\textbf{Armed}\\motors short-braked, holding the robot on the ramp};
    \node[run]  (climb) at (5.6,0)     {\textbf{Climb}\\both motors at full power (255)};
    \node[run]  (dec)   at (5.6,-3.5)  {\textbf{Decelerate}\\power ramps linearly to zero over \qty{150}{ms}};
    \node[run]  (hold)  at (5.6,-6.9)  {\textbf{Hold}\\motors short-braked};
    \node[run]  (back)  at (0,-6.9)    {\textbf{Back up}\\power 110};
    \node[run]  (fwd)   at (11.2,-6.9) {\textbf{Drive forward}\\power 255};

    \draw[->,thick] (armed) -- node[lbl,above] {start switch\\on} (climb);

    \draw[->,thick] (climb.north east) ++(-0.5,0)
      .. controls +(0.6,1.3) and +(-0.6,1.3) .. ([xshift=0.5cm]climb.north west)
      node[lbl,pos=0.5,above] {neither condition met};

    \draw[->,thick] (climb) -- node[lbl,right,align=left]
      {\qty{3.15}{s} after the start\\\textbf{or} pitch below \ang{18} for \qty{40}{ms}\\(pitch tested after the first \qty{0.3}{s})} (dec);

    \draw[->,thick] (dec) -- node[lbl,right,pos=0.35] {power at zero} (hold);

    \draw[->,thick] (hold.15) to[bend left=28] node[lbl,above] {nose up more\\than \ang{6}} (fwd.165);
    \draw[->,thick] (fwd.195) to[bend left=28] node[lbl,below] {back within\\\ang{3}} (hold.345);

    \draw[->,thick] (hold.165) to[bend right=28] node[lbl,above] {nose down more\\than \ang{6}} (back.15);
    \draw[->,thick] (back.345) to[bend right=28] node[lbl,below] {back within\\\ang{3}} (hold.195);

    \begin{scope}[shift={(9.6,0.55)},font=\footnotesize]
      \node[draw,fill=black!6,minimum size=2.6mm,inner sep=0pt] (k1) at (0,0) {};
      \node[anchor=west] at (0.2,0) {waiting for the start};
      \node[draw,fill=blue!12,minimum size=2.6mm,inner sep=0pt] (k2) at (0,-0.5) {};
      \node[anchor=west] at (0.2,-0.5) {autonomous run};
    \end{scope}
  \end{tikzpicture}
  \caption{The autonomous run. Climb ends on time or on pitch, whichever comes first. In Hold the robot reacts to pitch alone.}
  \label{fig:states}
\end{figure}

\subsection{Tuning}
\label{sec:tuning}

We empirically found that a climb takes about \qty{3.3}{s}, so we set \qty{3150}{ms} of full-power climb and a \qty{150}{ms} deceleration ramp. We also built a trim tool for straight-line driving that works as a binary search: each press shifts power \qty{16}{\percent} toward the lagging side, and the step halves every time the drift reverses. The side wall keeps the robot on a straight line either way, but a robot that pulls toward the wall spends power pushing itself into it, which only slows the climb. Trimming the motors to drive straight puts that power back into climbing.

At first, Hold corrected at full power, 255, in both directions. The robot then oscillated on the mesa: it threw itself forward, overshot, threw itself back, overshot, and so on. We lowered the back-up power to 110 to stop this, and kept forward correction at full power.

From our average climb speed of \qty{0.46}{m/s}, stopping takes $v^2/(2a)$: \qty{11}{cm} at \qty{1}{m/s^2} and \qty{4}{cm} at \qty{3}{m/s^2}, both inside the \qty{30}{cm} mesa. We never ran the robot without the Decelerate stage, so we cannot say whether overshooting the mesa was a real risk.
